\subsection{Flavor and the phase grammar}

The hard group-theoretic statement is
\[
G_{\mathrm{SM}}^{\mathrm{glob}}
=
S(U(3)\times U(2))
\cong \frac{SU(3)\times SU(2)\times U(1)_Y}{\mathbb Z_6}.
\]
The separate, still dynamical question is whether the observed CKM, PMNS, and Majorana phases actually realize the associated discrete grammar.

\begin{conjecture}[Phase grammar]
The global quotient data suggest a distinguished six-point phase lattice
\[
\delta \in \left\{
0,\frac{\pi}{3},\frac{2\pi}{3},\pi,\frac{4\pi}{3},\frac{5\pi}{3}
\right\}\mod 2\pi.
\]
With CP acting by inversion,
\[
\delta \mapsto -\delta \mod 2\pi,
\]
the resulting phase grammar is naturally
\[
\mathbb Z_6\rtimes \mathbb Z_2.
\]
\end{conjecture}

In v2.9 the quotient theorem should therefore be stated hard, while the physical realization of the phase grammar should remain explicitly conjectural.

\section{Conclusion}

The rigid core of v2.9 is not a number but a theorem:
\[
\boxed{
\text{admissibility on the orientable double cover and the seam involution}
}
\]
\[
\boxed{
\text{force the minimal carrier }E_3\oplus E_2.
}
\]
From this follow the global Standard Model group, one chiral family, the carrier norm $\tr_EY^2=5/6$, and the general abelian index formula. Only on the canonical occupancy branch does the same rigid package produce
\[
\Oadm = N_{\mathrm{fam}}\dim S^+ = 3\times 16 = 48.
\]
The center is therefore the minimal carrier. The number $48$ is the first large trace of that carrier packet in an admissible vacuum, not the foundation of the architecture.

\appendix

\section{v2.8 to v2.9 Dictionary}

\begin{table}[H]
\centering
\small
\begin{tabularx}{\linewidth}{@{}lX@{}}
\toprule
v2.8 language & v2.9 carrier reading \\
\midrule
$g=5$ & $\gamma=\frac56$ together with $B=\frac32$ on the minimal carrier packet \\
$48$ & $\Oadm=N_{\mathrm{fam}}\dim S^+$ on the canonical occupancy branch \\
$\bone=\frac{41}{10}$ & $\bone(N_{\mathrm{fam}},N_H)=\frac43 N_{\mathrm{fam}}+\frac1{10}N_H$ evaluated at $(3,1)$ \\
$\delta_2=\frac54\,\deltatop^2$ & $\delta_2 = B\gamma\,\Oadm^2\cthree^8$ with $B\gamma=\frac54$ \\
$\gamma(0)=\frac56$ in E$_8$ language & $\gamma=\tr_EY^2$ as a carrier invariant \\
\bottomrule
\end{tabularx}
\caption{Compact dictionary from the v2.8 closure language to the v2.9 carrier language.}
\end{table}

\section{Family-Index Bridge}

The v2.8 block relation
\[
k_B=\frac32\,I_1(B)
\]
already shows that family multiplicity is tied to the three distinguished cycles of the orientable double cover. Writing
\[
k_B=\frac12\sum_{i=1}^{3} I_{\mathrm{fam}}(C_i)\,I_1(B)
\]
makes the same factorization explicit and immediately yields
\[
\sum_{i=1}^{3} I_{\mathrm{fam}}(C_i)=3.
\]
The equal-one assignment is the minimal symmetric branch and is the natural source of $N_{\mathrm{fam}}=3$.

\section{Occupancy Transgression}

The canonical bridge is
\[
\Oadm = N_{\mathrm{fam}}\dim S^+ = 3\cdot 16 = 48.
\]
With $\cthree=1/(8\pi)$ this gives
\[
\Oadm\,\cthree^4
= 48\left(\frac{1}{8\pi}\right)^4
= \frac{3}{256\pi^4}
= \deltatop.
\]
This is the precise sense in which the topological surcharge is re-read as admissible occupancy transgression.

\section{Seam Descent Normalization}

The seam descent logic is intentionally separated into two pieces:
\begin{enumerate}[leftmargin=1.5em, topsep=3pt, itemsep=2pt]
    \item the spin-lifted seam loop closes only after $4\pi$, so $\widetilde U_\Gamma(2\pi)=-1$ and $\widetilde U_\Gamma(4\pi)=1$;
    \item if the admissible carrier modulus is identified with the axionic angle variable, then $a=\theta/(4\pi)$ and therefore $\cthree=1/(8\pi)$.
\end{enumerate}
The first step is geometric. The second is the explicit carrier-to-axion identification.

\section{Relative \texorpdfstring{$a_2$}{a2} Sketch}

The relative Einstein target rests on the standard four-dimensional heat-kernel structure
\[
a_2(D^2)=\frac{1}{16\pi^2}\int d^4x\sqrt{-g}\,
\tr\!\left(\mathcal E+\frac16 R\right).
\]
For Dirac-type curvature coupling $\mathcal E\sim R/4+\cdots$, so curvature terms naturally carry a denominator of the form $192\pi^2$. The open step is not whether such a denominator is plausible, but which relative subtraction and sector trace produce the final canonical coefficient
\[
\frac{\chi^2}{192\pi^2}\,\tr_{\Hch}\mathbf 1.
\]

\section{Numerical Branch Table}

\begin{table}[H]
\centering
\small
\begin{tabularx}{\linewidth}{@{}>{\raggedright\arraybackslash}Xcccc@{}}
\toprule
Branch & $N_{\mathrm{fam}}$ & $N_H$ & $\Oadm$ & $\alpha^{-1}(0)$ \\
\midrule
canonical carrier branch & $3$ & $1$ & $48$ & \num{137.03599921615843} \\
boundary multiplicity deformation & $2$ & $1$ & $32$ & \num{156.094045457} \\
boundary multiplicity deformation & $4$ & $1$ & $64$ & \num{124.835054800} \\
scalar-content deformation & $3$ & $2$ & $48$ & \num{135.945991742} \\
external truncation (no $\nu^c$) & $3$ & $1$ & $45$ & \num{137.033841585} \\
\bottomrule
\end{tabularx}
\caption{Representative branches in the carrier CFE. The ``no $\nu^c$'' row is not a normal carrier branch but an external truncation of the hard family packet.}
\end{table}

\end{document}
